\documentclass[11pt]{article}

\usepackage{acl}
\usepackage{times}
\usepackage{latexsym}
\usepackage{graphicx}
\usepackage{booktabs}
\usepackage{multirow}
\usepackage{amsmath}
\usepackage{amssymb}
\usepackage{url}
\usepackage{microtype}
\usepackage{xcolor}
\usepackage{pifont}
\newcommand{\cmark}{\ding{51}}
\newcommand{\xmark}{\ding{55}}

% -------------------------------------------------------
% Title and Authors
% -------------------------------------------------------
\title{TempoRAG-KG: Temporal Knowledge Graph-Augmented\\
Retrieval for Multi-Hop Question Answering}

\author{
  Supanut Kompayak \\
  st126055@ait.asia \\
  \textit{Asian Institute of Technology}
  \And
  Aphisit Jaemyaem \\
  st126130@ait.asia \\
  \textit{Asian Institute of Technology}
  \AND
  Dechathon Niamsa-ard \\
  st126235@ait.asia \\
  \textit{Asian Institute of Technology}
  \And
  Kaung Hein Htet \\
  st126477@ait.asia \\
  \textit{Asian Institute of Technology}
}

\begin{document}
\maketitle

% -------------------------------------------------------
% ABSTRACT
% -------------------------------------------------------
\begin{abstract}
Knowledge Graph-augmented Retrieval (KG-RAG) systems have demonstrated
strong performance on multi-hop question answering by structuring retrieved
facts as interconnected triples. However, existing systems---including the
recent state-of-the-art KG$^2$RAG---assume a static snapshot of knowledge
and fail to account for \textit{temporal validity}: the reality that facts
are only true within specific time windows. In this paper, we present
\textbf{TempoRAG-KG}, a temporally-aware KG-RAG framework that extends
KG$^2$RAG with two contributions: (1) explicit $[\mathtt{valid\_from},
\mathtt{valid\_to}]$ validity intervals attached to every KG triple,
enabling deterministic temporal filtering during retrieval, and (2) a
Graph-of-Graph (GoG)-based fill-in mechanism for handling incomplete KG
facts. We first conduct a systematic analysis of temporal failure modes in
KG-RAG, categorizing failures into four types and measuring their
distribution across HotpotQA and MuSiQue. Our exploratory data analysis
reveals that 35.0\% of HotpotQA and 38.3\% of MuSiQue development
questions contain temporal expressions, with prevalence increasing from
30.2\% in 2-hop to 46.7\% in 4-hop questions---suggesting temporal validity
becomes increasingly critical as reasoning complexity grows. We hypothesize
that TempoRAG-KG will improve F1 over the KG$^2$RAG baseline, with gains
concentrated in temporally-sensitive question subsets.
\end{abstract}

% -------------------------------------------------------
% 1. INTRODUCTION
% -------------------------------------------------------
\section{Introduction}

% Paragraph 1 — State of the world
Retrieval-Augmented Generation (RAG) has emerged as the dominant paradigm
for grounding large language model (LLM) responses in external factual
knowledge \cite{lewis2020retrieval}. Recent advances have moved beyond
flat vector retrieval toward structured graph-based approaches, where
Knowledge Graphs (KGs) provide fact-level relationships between retrieved
chunks. The state-of-the-art KG$^2$RAG \cite{zhu2025kg2rag} demonstrates
that organizing retrieved evidence through KG-guided expansion and Maximum
Spanning Tree context organization yields substantial improvements over
prior RAG baselines on multi-hop question answering benchmarks such as
HotpotQA \cite{yang2018hotpotqa}.

% Paragraph 2 — The big BUT
However, a fundamental limitation remains unaddressed: \textbf{all existing
KG-RAG systems, including KG$^2$RAG, assume a static knowledge graph in
which every triple is treated as universally and permanently valid.} In
reality, facts are only true within specific time periods---Tim Cook
\textit{became} CEO of Apple in 2011, implying Steve Jobs held that role
before him. When a user asks \textit{``Who was the CEO of Apple in 2005?''},
a system without temporal awareness retrieves both facts simultaneously and
either hallucinates or selects arbitrarily. We term this class of errors
\textit{temporal failure modes}, illustrated in Figure~\ref{fig:motivation}.
Despite the prevalence of such queries---our exploratory data analysis (EDA)
reveals that 35.0\% of HotpotQA and 38.3\% of MuSiQue development questions
contain temporal expressions---no existing KG-RAG framework systematically
addresses temporal validity at the triple level on general-domain unstructured
text.

% Paragraph 3 — Therefore, we did
To address this gap, we propose \textbf{TempoRAG-KG}, a framework that
extends KG$^2$RAG with two targeted components. First, we attach explicit
$[\mathtt{valid\_from}, \mathtt{valid\_to}]$ validity intervals to every
extracted triple during offline KG construction, and apply deterministic
temporal filtering prior to graph expansion. Second, we integrate a
GoG-based fill-in mechanism \cite{he2024gog} to recover answers when the
constructed KG is incomplete. Crucially, we precede system development
with a \textit{systematic diagnostic study} that categorizes temporal
failure modes into four types and quantifies their distribution across
two established multi-hop QA benchmarks.

% Paragraph 4 — Anticipated findings
We anticipate several key findings from our experiments. First, stale-fact
retrieval will constitute the majority of temporal errors in KG$^2$RAG,
as explicit year mention (17.8\% of HotpotQA) frequently surfaces
conflicting facts from different eras. Second, temporal prevalence will
escalate with reasoning complexity---from 30.2\% in 2-hop MuSiQue questions
to 46.7\% in 4-hop---suggesting that temporal awareness is disproportionately
important for complex multi-hop chains. Third, LLM-based temporal extraction
will achieve high precision on explicit date expressions but degrade on
implicit and conflicting temporal references, identifying a clear bottleneck
for future work.

% Paragraph 5 — Contributions
The contributions of this work are:
\begin{enumerate}
  \item \textbf{Temporal Failure Taxonomy}: A systematic categorization of
        four temporal failure modes in KG-RAG systems, with empirical
        distribution analysis across HotpotQA and MuSiQue.
  \item \textbf{TempoRAG-KG Framework}: A temporally-aware KG-RAG pipeline
        that attaches explicit validity intervals per triple and applies
        temporal filtering prior to graph expansion, extending
        KG$^2$RAG \cite{zhu2025kg2rag}.
  \item \textbf{Temporal Extraction Study}: An empirical evaluation of
        LLM-based temporal validity extraction from unstructured text,
        characterizing accuracy across three linguistic pattern types.
\end{enumerate}

% -------------------------------------------------------
% Figure 1 — Motivating example (in Introduction)
% -------------------------------------------------------
\begin{figure}[t]
  \centering
  \fbox{\parbox{0.45\textwidth}{
    \centering
    \vspace{1.2cm}
    \textbf{[Figure 1: Motivating Example]}\\[0.5em]
    \textbf{Query:} ``Who was the CEO of Apple in 2005?''\\[0.4em]
    \textbf{KG$^2$RAG retrieves:}\\
    (Steve Jobs, is\_CEO\_of, Apple) \\
    (Tim Cook, is\_CEO\_of, Apple) \\
    $\rightarrow$ \textcolor{red}{Conflicting — hallucinates}\\[0.4em]
    \textbf{TempoRAG-KG retrieves:}\\
    (Steve Jobs, is\_CEO\_of, Apple, \\ \quad valid\_from=1997, valid\_to=2011) \\
    Filter: $1997 \leq 2005 \leq 2011$ \checkmark \\
    $\rightarrow$ \textcolor{blue}{Correct: Steve Jobs}
    \vspace{1.2cm}
  }}
  \caption{Motivating example of a temporal failure mode in KG$^2$RAG
  (Type 2: Conflicting Facts). TempoRAG-KG resolves this by filtering
  edges with explicit validity intervals before retrieval.}
  \label{fig:motivation}
\end{figure}

% -------------------------------------------------------
% 2. RELATED WORK
% -------------------------------------------------------
\section{Related Work}

\subsection{KG-Augmented Retrieval-Augmented Generation}

Early RAG systems \cite{lewis2020retrieval} rely on dense passage retrieval
to retrieve isolated text chunks, ignoring structural relationships between
facts. Graph-based RAG approaches address this by incorporating KG structure
into retrieval. KG$^2$RAG \cite{zhu2025kg2rag} proposes KG-guided chunk
expansion using seed retrieval followed by a KG-based chunk organization
process via Maximum Spanning Tree, achieving strong performance on
HotpotQA. GEAR \cite{ma2025gear} introduces efficient beam-search graph
traversal that replaces expensive LLM calls at each hop with cosine
similarity scoring, substantially reducing inference cost. K-RagRec
\cite{wang2025kragrec} applies KG-guided RAG to recommendation, achieving
up to 93\% hallucination reduction. Despite these advances,
\textbf{none of these systems attach temporal validity metadata to KG
triples}, treating all retrieved facts as equally and permanently valid.

\subsection{Temporal Knowledge Graph Reasoning}

Temporal KG reasoning has been studied extensively over pre-built structured
KGs such as Wikidata and ICEWS. TempAgent \cite{tempagent2025} introduces
temporal filtering tools for an LLM agent, but requires a pre-existing
structured temporal KG and cannot operate over general-domain unstructured
text. EvoReasoner \cite{evoreasonerevokg2025} proposes a unified framework
combining multi-hop temporal reasoning with incremental KG evolution from
documents via confidence-based contradiction resolution and temporal trend
tracking---the closest existing work to TempoRAG-KG. However, EvoReasoner
tracks temporal evolution \textit{implicitly} through confidence scores
rather than explicit $[\mathtt{valid\_from}, \mathtt{valid\_to}]$ intervals,
which prevents deterministic point-in-time querying. Furthermore,
EvoReasoner is evaluated on structured TKGQA benchmarks (MultiTQ,
CronQuestions) rather than general-domain multi-hop QA datasets.
\textbf{Our work is the first to combine explicit triple-level temporal
validity tagging with general-domain KG-RAG evaluated on HotpotQA and
MuSiQue.}

\subsection{KG Construction from Unstructured Text}

Automatic KG construction enables RAG systems to operate without pre-built
KGs. HRGraph \cite{hrgraph2024} and STINMatch \cite{stinmatch2023} use LLM
prompting to extract subject-relation-object triples from raw text, while
TOBUGraph \cite{tobugraph2024} demonstrates production-scale construction.
GoG \cite{he2024gog} addresses KG incompleteness by allowing LLMs to fill
in missing facts from parametric knowledge when graph traversal reaches
dead ends. However, \textbf{none of these approaches extract or store
temporal validity intervals alongside triples}, losing the temporal context
necessary for time-sensitive queries. TempoRAG-KG extends LLM-based
extraction to additionally infer $\mathtt{valid\_from}$ and
$\mathtt{valid\_to}$ for each triple.

\paragraph{Gap Summary.}
No existing system simultaneously (1) automatically constructs a KG from
general-domain unstructured text, (2) attaches explicit per-triple temporal
validity intervals, and (3) evaluates on general-domain multi-hop QA
benchmarks. TempoRAG-KG fills this gap.

% -------------------------------------------------------
% 3. RESEARCH QUESTIONS
% -------------------------------------------------------
\section{Research Questions}

This work is organized around three research questions:

\begin{description}
  \item[\textbf{RQ1}] \textit{What types of temporal failures occur in
    KG-RAG systems when answering time-sensitive multi-hop questions,
    and how are they distributed across HotpotQA and MuSiQue?}\\
    \textbf{IV:} Question type (temporal vs.\ non-temporal), hop count.
    \textbf{DV:} Failure type distribution, failure rate per category.\\
    \textbf{H1:} Stale-fact retrieval (Type 1) will be the most frequent
    failure mode, and failure rate will increase with hop count in MuSiQue.

  \item[\textbf{RQ2}] \textit{To what extent does explicit triple-level
    temporal validity tagging improve retrieval precision and answer F1
    on temporally-sensitive questions compared to the KG$^2$RAG baseline?}\\
    \textbf{IV:} Presence of temporal tagging and filtering.
    \textbf{DV:} F1 score, Exact Match (EM), retrieval precision on
    temporal vs.\ non-temporal question subsets.\\
    \textbf{H2:} Temporal filtering will improve F1 on the temporal subset
    without degrading performance on non-temporal questions.

  \item[\textbf{RQ3}] \textit{How accurately can LLMs extract temporal
    validity intervals from unstructured text, and what linguistic patterns
    most affect extraction quality?}\\
    \textbf{IV:} Temporal expression type (explicit year, implicit/relative,
    conflicting).
    \textbf{DV:} Extraction precision, extraction recall, F1.\\
    \textbf{H3:} Extraction precision will be highest for explicit year
    mentions ($>$80\%) and lowest for conflicting date expressions.
\end{description}

% -------------------------------------------------------
% 4. METHODOLOGY
% -------------------------------------------------------
\section{Methodology}

% -------------------------------------------------------
% Figure 2 — System Architecture
% -------------------------------------------------------
\begin{figure*}[t]
  \centering
  \fbox{\parbox{0.95\textwidth}{
    \centering
    \vspace{2cm}
    \textbf{[Figure 2: TempoRAG-KG System Architecture]}\\[0.5em]
    \textit{Offline pipeline (top):}
    Documents $\rightarrow$ Chunking (512 tok, 100 overlap)
    $\rightarrow$ \textcolor{blue}{LLM Triple Extraction + valid\_from/valid\_to tagging}
    $\rightarrow$ \textcolor{blue}{Neo4j KG (temporal edges)} + FAISS Index\\[0.4em]
    \textit{Online pipeline (bottom):}
    Query $\rightarrow$ \textcolor{blue}{Temporal year detection ($y_q$)}
    $\rightarrow$ Seed retrieval (FAISS top-5)
    $\rightarrow$ \textcolor{blue}{Temporal filter (Eq.~1)}
    $\rightarrow$ GEAR beam search ($b$=3, $d$=2)
    $\rightarrow$ \textcolor{blue}{GoG fill-in (if dead end)}
    $\rightarrow$ MST context organization
    $\rightarrow$ LLM $\rightarrow$ Answer\\[0.3em]
    \small\textit{Blue = TempoRAG-KG additions; Gray = KG$^2$RAG backbone}
    \vspace{2cm}
  }}
  \caption{Overview of the TempoRAG-KG framework. Components in
  \textcolor{blue}{blue} are our additions to the KG$^2$RAG backbone.
  The offline pipeline constructs a temporally-annotated KG stored in
  Neo4j; the online pipeline applies temporal filtering before GEAR
  beam search and GoG fill-in for incomplete facts.}
  \label{fig:architecture}
\end{figure*}

\subsection{Datasets}

\paragraph{HotpotQA.}
We use the distractor development set of HotpotQA \cite{yang2018hotpotqa},
comprising 7,405 two-hop questions drawn from Wikipedia with ten candidate
paragraphs each. Our EDA reveals that \textbf{35.0\%} of questions
(2,593/7,405) contain temporal expressions. Bridge-type questions exhibit
higher temporal prevalence (39\%) than comparison-type questions (21\%),
consistent with the expectation that multi-hop reasoning more often
requires temporal grounding. Table~\ref{tab:eda} summarizes the
distribution by temporal pattern type.

\paragraph{MuSiQue.}
MuSiQue \cite{trivedi2022musique} provides 2,417 development questions
with 2--4 reasoning hops, designed to preclude shortcut reasoning.
\textbf{38.3\%} of questions (926/2,417) contain temporal expressions.
Crucially, temporal prevalence increases monotonically with hop count:
30.2\% for 2-hop, 47.2\% for 3-hop, and 46.7\% for 4-hop questions,
empirically motivating temporal-aware retrieval for complex reasoning chains.

\paragraph{Temporal Extraction Gold Standard.}
For RQ3, we manually annotate 200 passages sampled from the HotpotQA
training set, recording ground-truth $[\mathtt{valid\_from},
\mathtt{valid\_to}]$ for all extractable triples. Passages are stratified
across three temporal expression types: explicit year mention (e.g.,
``became CEO in 2011''), implicit/relative (e.g., ``during his tenure''),
and conflicting (multiple dates for the same relation).

% EDA Table
\begin{table}[t]
  \centering
  \small
  \begin{tabular}{lrr}
    \toprule
    \textbf{Temporal Pattern} & \textbf{HotpotQA} & \textbf{MuSiQue} \\
    \midrule
    Year mention        & 17.8\% & 6.7\%  \\
    When/what year      & 7.3\%  & 20.6\% \\
    Before/after        & 4.2\%  & 7.1\%  \\
    First/last/earliest & 7.7\%  & 7.1\%  \\
    Time period         & 2.3\%  & 3.9\%  \\
    Age/duration        & 0.7\%  & 1.8\%  \\
    \midrule
    \textbf{Any temporal} & \textbf{35.0\%} & \textbf{38.3\%} \\
    \bottomrule
  \end{tabular}
  \caption{Temporal expression prevalence in HotpotQA and MuSiQue
  development sets. Questions may match multiple patterns.}
  \label{tab:eda}
\end{table}

\subsection{Data Preprocessing}

\paragraph{Chunking.}
Each document is segmented into overlapping chunks of 512 tokens with a
100-token overlap, following KG$^2$RAG \cite{zhu2025kg2rag}. Each chunk
retains a reference to its source document.

\paragraph{Triple Extraction with Temporal Annotation.}
We prompt GPT-4 to extract structured triples from each chunk in the
format: $(\mathtt{head}, \mathtt{relation}, \mathtt{tail},
\mathtt{valid\_from}, \mathtt{valid\_to})$. The prompt instructs the model
to infer temporal validity from contextual cues (explicit years, relative
temporal expressions, succession language). If no temporal information is
available, fields are set to \texttt{null}. The full prompt is in
Appendix~\ref{app:prompt}.

\paragraph{KG Storage and Indexing.}
Extracted triples are stored in Neo4j with $\mathtt{valid\_from}$ and
$\mathtt{valid\_to}$ as edge properties. Entity names are embedded using
\texttt{all-MiniLM-L6-v2} \cite{reimers2019sentencebert} and stored in a
FAISS index for efficient similarity search.

\subsection{TempoRAG-KG Pipeline}

The online pipeline proceeds in six steps, illustrated in
Figure~\ref{fig:architecture}:

\textbf{(1) Temporal Year Detection.}
A regex-based classifier identifies whether the input query contains a
temporal expression and extracts the target year $y_q$.

\textbf{(2) Seed Retrieval.}
The query is embedded and the top-5 most similar chunks are retrieved
from the FAISS chunk index, following KG$^2$RAG.

\textbf{(3) Temporal Filtering.}
For each candidate edge $e$ with properties $v_f$ ($\mathtt{valid\_from}$)
and $v_t$ ($\mathtt{valid\_to}$), we retain $e$ if and only if:
\begin{equation}
  (v_f = \texttt{null} \;\lor\; v_f \leq y_q) \;\land\;
  (v_t = \texttt{null} \;\lor\; v_t \geq y_q)
  \label{eq:filter}
\end{equation}
Edges without temporal metadata are retained conservatively, ensuring
that extraction errors degrade gracefully.

\textbf{(4) GEAR Beam Search.}
Graph expansion proceeds via beam search \cite{ma2025gear} with beam
size $b=3$ and maximum depth $d=2$. Candidate neighbors are scored by
cosine similarity to the current subgoal embedding, avoiding LLM calls
during traversal.

\textbf{(5) GoG Fill-in.}
When expansion reaches a dead end (no outgoing edges satisfy
Equation~\ref{eq:filter}), we invoke a GoG-style LLM fill-in
\cite{he2024gog}: the model supplies the missing fact from parametric
knowledge along with a confidence score. Fill-in facts with confidence
below 0.7 are discarded.

\textbf{(6) MST Context Organization and Generation.}
Retrieved subgraph facts are organized via Maximum Spanning Tree pruning
into coherent paragraphs, prepended to the generation prompt for
LLaMA-3-8B or GPT-4.

\subsection{Baselines}

We compare against four baselines:
\begin{itemize}
  \item \textbf{Vanilla RAG}: Dense retrieval with GPT-4; no graph.
  \item \textbf{GraphRAG}: Microsoft community-summary RAG \cite{edge2024graphrag}.
  \item \textbf{KG$^2$RAG}: Reproduced NAACL 2025 baseline \cite{zhu2025kg2rag}.
  \item \textbf{TempoRAG-KG (ours)}: Full system with temporal filter + GoG.
\end{itemize}

\subsection{Ablation Design}

To isolate each component's contribution (RQ2), we evaluate five conditions:

\begin{table}[t]
  \centering
  \small
  \begin{tabular}{lccc}
    \toprule
    \textbf{Condition} & \textbf{Temporal} & \textbf{GoG} & \textbf{Graph} \\
    \midrule
    Vanilla RAG            & \xmark & \xmark & \xmark \\
    KG$^2$RAG              & \xmark & \xmark & \cmark \\
    +Temporal only         & \cmark & \xmark & \cmark \\
    +GoG only              & \xmark & \cmark & \cmark \\
    \textbf{Full (ours)}   & \cmark & \cmark & \cmark \\
    \bottomrule
  \end{tabular}
  \caption{Ablation conditions. \cmark\ = component enabled;
  \xmark\ = disabled.}
  \label{tab:ablation}
\end{table}

\subsection{Evaluation Metrics}

We report:
\begin{itemize}
  \item \textbf{F1 and EM}: Token-level F1 and Exact Match, computed
    overall and separately for temporal and non-temporal subsets.
  \item \textbf{Retrieval Precision@5}: Fraction of retrieved chunks
    containing the gold supporting facts.
  \item \textbf{Temporal Extraction F1}: Precision and recall against
    gold $[\mathtt{valid\_from}, \mathtt{valid\_to}]$ annotations
    with $\pm$1 year tolerance (RQ3).
  \item \textbf{Inference Time}: Average seconds per query.
\end{itemize}

% -------------------------------------------------------
% Figure 3 — Temporal Failure Taxonomy
% -------------------------------------------------------
\begin{figure}[t]
  \centering
  \fbox{\parbox{0.45\textwidth}{
    \centering
    \vspace{1.2cm}
    \textbf{[Figure 3: Temporal Failure Taxonomy]}\\[0.5em]
    2$\times$2 grid:\\[0.3em]
    \textbf{Type 1 — Stale Fact}\\
    Expired fact retrieved as current\\
    \textit{``Steve Jobs'' as Apple CEO now}\\[0.3em]
    \textbf{Type 2 — Conflicting Facts}\\
    Multiple contradictory facts returned\\
    \textit{Both Jobs and Cook returned}\\[0.3em]
    \textbf{Type 3 — Missing Temporal Context}\\
    Correct entity, unknown validity window\\
    \textit{Cook is CEO, but which years?}\\[0.3em]
    \textbf{Type 4 — Temporal Hop Failure}\\
    Stale fact at any hop breaks the chain\\
    \textit{Hop 1 correct; hop 2 retrieves outdated}
    \vspace{1.2cm}
  }}
  \caption{Four types of temporal failure modes in KG-RAG, identified
  through qualitative analysis of KG$^2$RAG errors on the HotpotQA
  temporal subset. TempoRAG-KG directly addresses Types 1--3 via
  temporal filtering and Type 4 via GoG fill-in.}
  \label{fig:taxonomy}
\end{figure}

% -------------------------------------------------------
% 5. EXPECTED RESULTS
% -------------------------------------------------------
\section{Expected Results}

% Main Results Table
\begin{table*}[t]
  \centering
  \small
  \begin{tabular}{lcccccc}
    \toprule
    \multirow{2}{*}{\textbf{Model}} &
    \multicolumn{2}{c}{\textbf{Overall}} &
    \multicolumn{2}{c}{\textbf{Temporal Subset}} &
    \multicolumn{2}{c}{\textbf{Non-temporal Subset}} \\
    \cmidrule(lr){2-3}\cmidrule(lr){4-5}\cmidrule(lr){6-7}
    & \textbf{F1} & \textbf{EM}
    & \textbf{F1} & \textbf{EM}
    & \textbf{F1} & \textbf{EM} \\
    \midrule
    Vanilla RAG                     & -- & -- & -- & -- & -- & -- \\
    GraphRAG \cite{edge2024graphrag} & -- & -- & -- & -- & -- & -- \\
    KG$^2$RAG \cite{zhu2025kg2rag}  & -- & -- & -- & -- & -- & -- \\
    \midrule
    \textbf{TempoRAG-KG (ours)}     & \textbf{--} & \textbf{--}
                                    & \textbf{--} & \textbf{--}
                                    & \textbf{--} & \textbf{--} \\
    \bottomrule
  \end{tabular}
  \caption{Expected main results on HotpotQA distractor dev set.
  Temporal and non-temporal subsets contain 2,593 and 4,812 questions
  respectively. \textbf{Bold} = best result. Results to be filled
  after experiments.}
  \label{tab:main_results}
\end{table*}

% Ablation Scores Table
\begin{table}[t]
  \centering
  \small
  \begin{tabular}{lcc}
    \toprule
    \textbf{Condition} & \textbf{Overall F1} & \textbf{Temporal F1} \\
    \midrule
    Vanilla RAG                     & -- & -- \\
    KG$^2$RAG (backbone)            & -- & -- \\
    \quad +Temporal filter only     & -- & -- \\
    \quad +GoG fill-in only         & -- & -- \\
    \quad \textbf{Full model (ours)}& \textbf{--} & \textbf{--} \\
    \midrule
    $\Delta$ vs.\ KG$^2$RAG        & --  & -- \\
    \bottomrule
  \end{tabular}
  \caption{Ablation scores on HotpotQA.
  $\Delta$ = improvement over KG$^2$RAG backbone.
  Results to be filled after experiments.}
  \label{tab:ablation_scores}
\end{table}

% Extraction Table
\begin{table}[t]
  \centering
  \small
  \begin{tabular}{lccc}
    \toprule
    \textbf{Pattern Type} & \textbf{Prec.} & \textbf{Recall} & \textbf{F1} \\
    \midrule
    Explicit year mention & -- & -- & -- \\
    Implicit / relative   & -- & -- & -- \\
    Conflicting dates     & -- & -- & -- \\
    \midrule
    \textbf{Overall}      & \textbf{--} & \textbf{--} & \textbf{--} \\
    \bottomrule
  \end{tabular}
  \caption{Expected temporal extraction accuracy on 200 manually
  annotated HotpotQA passages (RQ3), with $\pm$1 year tolerance.
  Results to be filled after annotation.}
  \label{tab:extraction}
\end{table}

\paragraph{RQ1: Temporal Failure Distribution.}
We expect stale-fact retrieval (Type 1) to constitute the majority of
temporal errors in KG$^2$RAG, consistent with year-mention being the most
frequent pattern (17.8\%). Type 4 (temporal hop failure) is expected to be
rare in HotpotQA (2-hop only) but more frequent in MuSiQue 4-hop questions.
This taxonomy, to our knowledge, is the first systematic categorization of
temporal failures in general-domain KG-RAG.

\paragraph{RQ2: Impact of Temporal Tagging.}
We expect the full model (Table~\ref{tab:ablation_scores}) to show gains
over KG$^2$RAG concentrated in the temporal subset, while non-temporal F1
remains stable---confirming that temporal filtering does not harm general
retrieval. Temporal and GoG components are expected to contribute
independently, as they target distinct failure modes.

\paragraph{RQ3: Temporal Extraction Accuracy.}
We expect precision $>$80\% for explicit year-mention expressions but
substantially lower recall for implicit/relative expressions (e.g.,
``during his tenure'') where temporal boundaries must be inferred.
Conflicting date expressions are expected to pose the greatest challenge,
directly characterizing the practical ceiling of TempoRAG-KG's filtering
quality.

% -------------------------------------------------------
% 6. DISCUSSION
% -------------------------------------------------------
\section{Discussion}

\paragraph{Design Implications.}
If our hypotheses are confirmed, the results will carry a broad implication
for KG-RAG system design: \textit{temporal metadata should be treated as a
first-class property of KG edges, not an optional annotation.} Systems that
attach validity intervals at construction time gain the ability to filter
stale facts deterministically---a capability that cannot be retrofitted to
static KGs without re-extraction. This suggests that future KG construction
pipelines should routinely include temporal extraction as a component
alongside entity and relation extraction.

\paragraph{Scope of Temporal Validity.}
TempoRAG-KG addresses point-in-time validity (``was this true at year
$y_q$?'') but does not model temporal trends, rates of change, or
predictive validity (``will this still be true next year?''). For rapidly
evolving domains such as financial news or sports statistics, a dynamic
update mechanism such as EvoKG \cite{evoreasonerevokg2025} may be
preferable. Our approach is designed for document corpora where knowledge
changes on a yearly rather than daily timescale.

\paragraph{Limitations.}
The primary limitation is extraction quality on implicit temporal
expressions. If extraction F1 on implicit patterns falls substantially,
the temporal filter introduces false negatives, discarding valid facts.
Our conservative fallback (retaining triples with \texttt{null} validity)
mitigates but does not eliminate this risk. A second limitation is that
our gold standard annotation (200 passages) is relatively small; a
larger study would provide more reliable estimates of extraction difficulty
by pattern type. Finally, all experiments use HotpotQA and MuSiQue, which
are English Wikipedia-based; generalization to other domains and languages
requires further investigation.

% -------------------------------------------------------
% 7. PROGRESS SUMMARY
% -------------------------------------------------------
\section{Progress Summary}

\paragraph{Completed.}
\begin{itemize}
  \item Literature review: 16 papers across four team members
    (4 papers each, per course requirements from ACL venues only).
  \item EDA of HotpotQA (7,405 questions: 35.0\% temporal) and
    MuSiQue (2,417 questions: 38.3\% temporal) using automated
    regex-based classifier with seven pattern categories.
  \item Temporal failure taxonomy: four types proposed based on
    qualitative analysis of KG-RAG failure modes.
  \item Analysis of closest competitor EvoReasoner \cite{evoreasonerevokg2025}
    with clear differentiation on three dimensions: explicit vs.\
    implicit temporal model, general-domain vs.\ structured KG,
    and benchmark choice.
\end{itemize}

\paragraph{In Progress.}
\begin{itemize}
  \item Neo4j schema design and setup with temporal edge properties.
  \item Temporal extraction prompt engineering and pilot evaluation
    on 20 HotpotQA passages.
  \item GitHub repository setup with README.\footnote{
    Repository: \url{https://github.com/[to-be-updated]/temporag-kg}}
\end{itemize}

\paragraph{Remaining Work.}
\begin{itemize}
  \item Reproduction of KG$^2$RAG baseline on HotpotQA distractor dev set.
  \item Manual annotation of 200 passages for RQ3 gold standard.
  \item Full pipeline implementation: temporal filter, GEAR beam search,
    GoG fill-in integration.
  \item Ablation experiments (Table~\ref{tab:ablation}).
  \item Streamlit demo with edge case scenarios.
\end{itemize}

\paragraph{Potential Concerns.}
The primary risk is temporal extraction quality on implicit expressions.
If extraction F1 on implicit patterns falls below 60\%, the temporal
filter may introduce false negatives. The conservative null-retention
strategy in Equation~\ref{eq:filter} ensures graceful degradation.
A secondary risk is KG$^2$RAG reproduction fidelity; we will verify
outputs against reported HotpotQA scores before running comparisons.

% -------------------------------------------------------
% 8. CONCLUSION
% -------------------------------------------------------
\section{Conclusion}

We have presented TempoRAG-KG, a temporally-aware Knowledge Graph-augmented
RAG framework that targets a systematic and underexplored failure mode in
existing KG-RAG systems: the inability to distinguish facts valid at
different points in time. By attaching explicit validity intervals to KG
triples and filtering graph expansion accordingly, TempoRAG-KG addresses
the 35\% of HotpotQA and 38\% of MuSiQue questions that require temporal
reasoning. Our temporal failure taxonomy provides the diagnostic foundation
needed to understand when and why temporal reasoning fails in
graph-augmented systems, and our planned extraction study characterizes
the practical accuracy ceiling of LLM-based temporal annotation. This
work contributes to the growing intersection of structured knowledge
retrieval and temporal reasoning, and opens a tractable path toward
knowledge systems that remain reliably accurate across time.

% -------------------------------------------------------
% REFERENCES
% -------------------------------------------------------
\bibliographystyle{acl_natbib}
\bibliography{references}

\appendix

% -------------------------------------------------------
% APPENDIX A — Temporal Extraction Prompt
% -------------------------------------------------------
\section{Temporal Extraction Prompt}
\label{app:prompt}

\begin{verbatim}
You are a knowledge extraction assistant.
Given the following text passage, extract all
factual triples. For each triple, also infer
temporal validity if mentioned or implied.

Format each triple as JSON:
{
  "head": "entity name",
  "relation": "relation type",
  "tail": "entity or value",
  "valid_from": year (int) or null,
  "valid_to":   year (int) or null,
  "confidence": float between 0 and 1
}

Rules:
- Only extract explicitly stated facts
- valid_from/valid_to: null if unknown
- Ongoing facts: set valid_to to null
- Conflicting dates: extract as separate triples

Text: {chunk_text}
\end{verbatim}

% -------------------------------------------------------
% APPENDIX B — EDA Methodology
% -------------------------------------------------------
\section{EDA Methodology}
\label{app:eda}

Temporal questions were identified using seven regex patterns applied
to the lowercased question string. A question was classified as temporal
if it matched at least one pattern. The patterns target:
(1) explicit year mentions,
(2) when/what-year phrases,
(3) before/after/since,
(4) time period expressions,
(5) first/last/earliest superlatives,
(6) age and duration phrases, and
(7) temporal relation words (predecessor, successor, contemporary).
The EDA script and full pattern definitions are available in the
project repository.

\end{document}
